% Shared macros for the main text and the appendices.
% Loaded by paper.tex after \usepackage[accepted]{mlsys2025}. Appendix-only macros belong
% in the appendix files; a macro both parts use belongs here.
% ---------------------------------------------------------------------------
% Packages. The MLSys style already loads times, fancyhdr, color, natbib, algorithm,
% algorithmic, eso-pic and forloop, and redefines \cite as \citep.
\usepackage{amsmath,amssymb,amsthm,mathtools,bm}
\usepackage{newtxmath}
\usepackage[scaled=.86,varqu]{zi4}
\usepackage[scaled=.9]{helvet}
\usepackage{booktabs,tabularx,array,multirow,longtable}
\usepackage{xcolor,tikz,pgfplots,pgfplotstable}
\usetikzlibrary{arrows.meta,positioning,calc,fit,backgrounds,decorations.pathreplacing,patterns}
\pgfplotsset{compat=1.18}
\usepgfplotslibrary{groupplots,fillbetween}
\usepackage{enumitem,etoolbox}
\usepackage{float}% [H]: the register's source table stays in Appendix F.1
\usepackage{xurl}
% ---------------------------------------------------------------------------
% Colours and figure styles: figures/style.tex (palette, series, axes, CSV look-ups).
% Shared figure style, loaded by macros.tex. Every figure in the paper and the
% research notes takes its colours, line styles, axis conventions and panel titles from this file,
% so a colour or a series means the same thing wherever it appears.
% ---------------------------------------------------------------------------
% Palette. Each colour has one meaning in every figure of the paper and the research notes.
% Method colours (three hues from Okabe and Ito's set and a grey):
%   cstock     ink             the stock engine's baseline: plain decoding (the target alone), the
%                              stock head, the target's head and head input h^t, reference lines;
%   ccert      blue            the certified head and its int8 copies (the proposed method);
%   cfallback  vermilion       fallback to the stock head, and decisions left open;
%   ccontext   grey            context and alternatives not pursued: other kernels, the FP8 and int4
%                              copies, dropped entries, envelope bands, and series with no
%                              paper-wide meaning (the domains of a panel, constructed head families).
% Arm colours, one per speculative serving arm, also used for the drafter it runs:
%   cmtp       green           MTP (the native MTP layer, MTP-4B) and its drift;
%   cdflash16  reddish purple  DFlash with 16-token blocks, and DFlash-4B wherever it drafts 16-token
%                              blocks (the captures, the drafter panel): h^d, Delta and its drift;
%   cdflash8   violet          DFlash with 8-token blocks.
% Plain decoding is the target alone and takes the stock ink. Series that differ only in the stock
% engine's configuration (the module kinds of the divergence figure, the lever stacks of plain
% decoding) stay ink and are told apart by marker and dash pattern. No hue doubles for another
% meaning: a drafter's head is drawn in its arm's colour, never in a generic draft colour.
% Checked with the dataviz skill's validate_palette.js (Machado, Oliveira and Fernandes 2009
% simulation at severity 1, OKLab distances x100, every pair compared, white surface). The five
% hues pass every check: OKLCH lightness 0.44-0.68, chroma at least 0.12, contrast at least
% 3.06:1 on white (reddish purple; the rest 3.87:1 or more); the closest pair under colour-vision
% deficiency is green and vermilion (8.2, protanopia), the closest in normal vision reddish purple
% and vermilion (16.4). The neutrals: ink is at least 13.8 from every hue under protanopia and
% deuteranopia (violet); the grey is 7.6 from reddish purple (protanopia) and 11.5 from blue in
% normal vision, so a grey series always has its own dash pattern and a direct label. The hues
% have similar lightness, so a greyscale print separates series only by line style or marker:
% every series has its own dash pattern or marker as well.
\definecolor{ink}{HTML}{1B3A57}
\definecolor{vpblue}{HTML}{0072B2}
\definecolor{vpvermilion}{HTML}{D55E00}
\definecolor{muted}{HTML}{6B7780}
\definecolor{vpgreen}{HTML}{1C6402}
\definecolor{vppurple}{HTML}{CC79A7}
\definecolor{vpviolet}{HTML}{6F1FA7}
\colorlet{cstock}{ink}
\colorlet{ccert}{vpblue}
\colorlet{cfallback}{vpvermilion}
\colorlet{ccontext}{muted}
\colorlet{cmtp}{vpgreen}
\colorlet{cdflash16}{vppurple}
\colorlet{cdflash8}{vpviolet}
% Fills are light tints of the hue they belong to; wash is the neutral one for shaded regions.
\colorlet{wash}{ink!6}
% ---------------------------------------------------------------------------
% Status of what a mark shows. Measured values carry filled markers, and empirical
% distributions carry sparse markers (4 to 8 per curve); a hollow marker is a measured subset its
% label or the caption names; values computed point by point from measurements carry the same markers with
% their own dash pattern and say so in their label; curves derived from a formula, a model or the
% weights are plain lines without markers; schematics have no numeric ticks and say "schematic" in
% the panel title, or once in the caption when every panel is one. Uncertainty is drawn where
% the data have it (whiskers or a band), or the caption gives the spread when it is smaller than a
% marker.
% A figure with several panels titles each one above its plot area: its letter in bold, a short
% title and the status (measured, derived, expected, schematic, to scale). A figure with one panel
% has no title inside the figure; its caption names it and its status. Every series is named by a
% direct label beside its line, mark or segment, at least 2 pt from any other line, mark or label,
% with a thin leader line where the free space is a little way off. A legend is used only where
% no such place exists, and the figure's comment says why; it then sits inside the plot in an
% empty region, covering no data, or above several panels it serves. Never a legend below an x
% label or under one panel only.
\newcommand{\paneltitle}[2]{\textbf{(#1)}~#2}
% ---------------------------------------------------------------------------
% Lines and marks. Data series are thick (0.8 pt); reference lines (the stock head, a bound of
% the text) semithick ink; leader lines, callouts and dimension arrows thin; very thick only for
% vectors drawn to scale. Marker sizes give the four shapes about the same area. Marker shape
% says what a series is: Hopper model a circle and conservative model a square wherever both models
% are drawn; plain decoding a circle, DFlash a diamond and MTP a triangle wherever arms or
% drafters are drawn.
\pgfplotsset{
  mark circle/.style={mark=*,mark size=1.6pt,mark options={solid}},
  mark square/.style={mark=square*,mark size=1.4pt,mark options={solid}},
  mark triangle/.style={mark=triangle*,mark size=2pt,mark options={solid}},
  mark diamond/.style={mark=diamond*,mark size=2pt,mark options={solid}},
  % a measured subset: the same shape, white inside, 0.6 pt outline
  mark hollow/.style={mark options={solid,fill=white,line width=0.6pt}},
  % Plot-area heights, so that panels in one row match: a full panel, and each of two panels
  % stacked on one x axis.
  panel full/.style={scale only axis,height=2.8cm},
  panel stacked/.style={scale only axis,height=1.75cm},
}
\tikzset{
  ref line/.style={semithick,draw=cstock},
  leader/.style={thin,draw=ink!70},
}
% ---------------------------------------------------------------------------
% Series. One style per quantity that recurs across figures.
\pgfplotsset{
  % The two error models of the stock kernel, the same pair everywhere: the Hopper model solid
  % with circles, the conservative model dashed with squares. Combine with a colour.
  model hopper/.style={solid,mark circle},
  model cons/.style={dashed,mark square},
  derived/.style={mark=none},
  % Copies of the head: the int8 copies are the method (blue); FP8 and int4 are alternatives (grey).
  copy int8 row/.style={ccert,thick},
  copy int8 g128/.style={ccert,thick,dashed},
  copy fp8/.style={ccontext,thick,dashdotted},
  copy int4/.style={ccontext,thick,densely dotted},
  % Arms and their drafters, the same wherever they appear (serving frontier, drift
  % distributions, drafter panel): plain decoding ink, solid, circles; DFlash with 8-token blocks
  % violet, solid, diamonds; DFlash with 16-token blocks reddish purple, dashed, diamonds; MTP
  % green, dash-dotted, triangles.
  arm plain/.style={cstock,thick,mark circle},
  arm dflash8/.style={cdflash8,thick,mark diamond},
  arm dflash16/.style={cdflash16,thick,dashed,mark diamond},
  arm mtp/.style={cmtp,thick,dashdotted,mark triangle},
}
% ---------------------------------------------------------------------------
% Text and axes. Three sizes, set in points so that a figure prints the same
% in the paper (MLSys: \footnotesize 9 pt) and in the research notes (article 10 pt: \footnotesize
% 8 pt): axis labels and panel titles 9 pt (\figlabel); sentences inside schematics and to-scale
% diagrams 8 pt (\figsmall, style note); tick labels, legends, keys, direct labels and values
% printed in a plot 7 pt (\scriptsize, style annot), the smallest size used. Text is black; the
% coloured line or mark beside a label carries the identity. Labels and titles in sentence case.
% Tick labels are plain numbers (100, 1k, 0.001), not powers of ten, whose exponents would fall
% below 7 pt. Figures are drawn at their final size: never scale a figure.
% Every figure file ends with a %, and so does every line that \inputs one: a space after the
% picture lets a full-width figure overflow its line, which pushes the caption's link anchor onto
% a line of its own and leaves a blank line (about 12 pt) between the figure and its caption.
% Tables use \tablefont (9 pt) in both documents, never smaller than 7 pt, natural width when they fit.
% \vpnum{digits}{expression} and \vppct{digits}{expression} print a computed value.
\newcommand{\figlabel}{\fontsize{9}{10}\selectfont}
\newcommand{\figsmall}{\fontsize{8}{9.5}\selectfont}
\newcommand{\tablefont}{\fontsize{9}{10}\selectfont}
\newcommand{\vpnum}[2]{\pgfmathparse{#2}\pgfmathprintnumber[fixed,precision=#1]{\pgfmathresult}}
\newcommand{\vppct}[2]{\pgfmathparse{#2}\pgfmathprintnumber[fixed,precision=#1,fixed zerofill]{\pgfmathresult}\%}
\tikzset{
  panel title/.style={font=\figlabel,text height=1.6ex,text depth=.4ex,inner sep=2pt},
  % An arrow means "passes to"; every arrow carries a label saying what passes or when.
  flow/.style={-{Stealth[length=1.6mm]},semithick,draw=ink},
  note/.style={font=\figsmall,text=black,align=center},
  annot/.style={font=\scriptsize,inner sep=1pt,text=black},
}
\pgfplotsset{every axis/.append style={
    tick label style={font=\scriptsize},label style={font=\figlabel},
    title style={panel title,yshift=-2pt},
    axis line style={ink!60},tick style={ink!60},major tick length=2.5pt,
    grid style={draw=ink!10},
    legend style={font=\scriptsize,draw=none,fill=white,fill opacity=.85,text opacity=1,
      inner sep=1.5pt,row sep=-1pt},
    legend cell align=left},
  % Keep only rows whose column #1 equals #2 (long-form CSVs).
  discard if not/.style 2 args={x filter/.append code={%
    \edef\tempa{\thisrow{#1}}\edef\tempb{#2}%
    \ifx\tempa\tempb\else\def\pgfmathresult{inf}\fi}},
}
% ---------------------------------------------------------------------------
% Values from committed CSVs. \vpreadcsv{\table}{path} reads a CSV once;
% \vplookup{\table}{key column}{key}{value column}{\result} sets \result to the value
% in the first row whose key column equals the key, and stops the build if none does.
% \vplookuptwo and \vplookupthree match two or three columns:
% \vplookupthree{\table}{column 1}{key 1}{column 2}{key 2}{column 3}{key 3}{value column}{\result}.
% Keys are expanded, so a key may be a macro. Figures use these instead of typing a number read
% from a file; values that exist only in JSON are copied to figures/data/json_values.csv by
% figures/data/make_figure_data.py (figures/data/README.md).
\makeatletter
\newif\ifvp@match
\newcommand{\vpreadcsv}[2]{\pgfplotstableread[col sep=comma]{#2}#1}
\newcommand{\vplookup}[5]{\vplookupthree{#1}{#2}{#3}{}{}{}{}{#4}{#5}}
\newcommand{\vplookuptwo}[7]{\vplookupthree{#1}{#2}{#3}{#4}{#5}{}{}{#6}{#7}}
\newcommand{\vplookupthree}[9]{%
  \global\let#9\relax
  \edef\vp@keya{#3}\edef\vp@keyb{#5}\edef\vp@keyc{#7}\def\vp@colb{#4}\def\vp@colc{#6}%
  \pgfplotstableforeachcolumnelement{#2}\of#1\as\vp@cell{%
    \ifx#9\relax\ifx\vp@cell\vp@keya
      \vp@matchtrue
      \ifx\vp@colb\@empty\else
        \pgfplotstablegetelem{\pgfplotstablerow}{#4}\of#1\let\vp@cellb\pgfplotsretval
        \ifx\vp@cellb\vp@keyb\else\vp@matchfalse\fi
      \fi
      \ifx\vp@colc\@empty\else
        \pgfplotstablegetelem{\pgfplotstablerow}{#6}\of#1\let\vp@cellc\pgfplotsretval
        \ifx\vp@cellc\vp@keyc\else\vp@matchfalse\fi
      \fi
      \ifvp@match\vp@take{#1}{#8}{#9}\fi
    \fi\fi}%
  \ifx#9\relax\PackageError{paper}{No row with #2 = \vp@keya\space in the table}{}\fi}
\newcommand{\vp@take}[3]{\pgfplotstablegetelem{\pgfplotstablerow}{#2}\of#1\global\let#3\pgfplotsretval}
% \vpsum{\table}{key column}{key}{value column}{\result}: the sum of the value column over
% the rows whose key column equals the key, printed as an integer.
\newcommand{\vpsum}[5]{%
  \gdef#5{0}\edef\vp@keya{#3}%
  \pgfplotstableforeachcolumnelement{#2}\of#1\as\vp@cell{%
    \ifx\vp@cell\vp@keya
      \pgfplotstablegetelem{\pgfplotstablerow}{#4}\of#1%
      \pgfmathtruncatemacro\vp@s{#5+\pgfplotsretval}\global\let#5\vp@s
    \fi}}
% \vpjson{key}{\result}: a value copied from JSON evidence into figures/data/json_values.csv.
\newcommand{\vpjson}[2]{%
  \ifdefined\vp@jsontable\else\vpreadcsv{\vp@jsontable}{figures/data/json_values.csv}\fi
  \vplookup{\vp@jsontable}{key}{#1}{value}{#2}}
\makeatother

\hypersetup{colorlinks=true,linkcolor=ink,citecolor=ink,urlcolor=ink}
% ---------------------------------------------------------------------------
% Lists and layout helpers.
\setlist{nosep,leftmargin=1.3em,topsep=.4ex,itemsep=.2ex}
\newcolumntype{P}[1]{>{\raggedright\arraybackslash}p{#1}}
\newcolumntype{Y}{>{\raggedright\arraybackslash}X}
% ---------------------------------------------------------------------------
% Mathematics.
\newcommand{\E}{\mathbb E}
\newcommand{\R}{\mathbb R}
\newcommand{\BF}{\mathbb B}
\newcommand{\ind}{\mathbf 1}
\DeclareMathOperator*{\argmax}{arg\,max}
\DeclareMathOperator{\softmax}{softmax}
\DeclareMathOperator{\TV}{TV}
\DeclareMathOperator{\KL}{KL}
\DeclareMathOperator{\fl}{fl}
\DeclareMathOperator{\diam}{diam}
\newcommand{\norm}[1]{\lVert #1\rVert}
\newcommand{\rnb}{\operatorname{rn}_{\mathrm{B}}}
\newcommand{\ulpb}{\operatorname{ulp}_{\mathrm{B}}}
% The ends of an interval, set as words so that hi_i does not read as h times i_i.
\newcommand{\lo}{\mathit{lo}}
\newcommand{\hi}{\mathit{hi}}
\newcommand{\code}[1]{\texttt{#1}}
\theoremstyle{plain}
\newtheorem{theorem}{Theorem}[section]
\newtheorem{proposition}[theorem]{Proposition}
\newtheorem{lemma}[theorem]{Lemma}
\theoremstyle{definition}
\newtheorem{definition}[theorem]{Definition}
% ---------------------------------------------------------------------------
% Names used throughout. The four references of a single decision (Definition
% def:contracts) and the two named models of the stock kernel's FP32 accumulation error
% (Table tab:assumptions). Always name the model beside a fallback share or a
% divergence class.
\newcommand{\Rreal}{\mbox{R-real}}
\newcommand{\Rbf}{\mbox{R-bf16}}
\newcommand{\Rstock}{\mbox{R-stock}}
\newcommand{\Srace}{\mbox{S-race}}
\newcommand{\modelcons}{conservative model}
\newcommand{\modelhopper}{Hopper model}
\newcommand{\gcons}{6.11\times10^{-4}}
\newcommand{\ghop}{1.193\times10^{-4}}
% Pins.
\newcommand{\sourcepin}{\texttt{bd66ce343e4f}}
\newcommand{\modelpin}{\texttt{851bf6e806ef}}
\newcommand{\repourl}{\url{https://github.com/SamMausberg/verified-progress}}
% The date of the build, in the paper's day-month-year style: the paper's preprint notice and the
% research notes' title page print it, so both documents carry the date they were built.
\newcommand{\builddate}{\number\day\space\ifcase\month\or January\or February\or March\or April\or May\or June\or July\or August\or September\or October\or November\or December\fi\space\number\year}
% ---------------------------------------------------------------------------
% Evidence status. A committed result cites its entry in the evidence register
% (appendix, \label{ev:<key>}): \ev{key} or \ev{key1,key2} prints a small tag, [E3] or [E3,E7], in
% script size so that it reads lighter than the text (a superscript would look like an exponent
% after a number). Work that was not done is stated as such in the text, without a marker.
\newcommand{\evref}[1]{\hyperref[ev:#1]{E\ref*{ev:#1}}}
\newcommand{\evloop}[1]{\evsep\evref{#1}\def\evsep{,}}
\DeclareRobustCommand{\ev}[1]{{\scriptsize\def\evsep{}[\forcsvlist{\evloop}{#1}]}}
% A figure that reads a committed data file: draw it if the file exists, otherwise a box naming
% the missing file, so that a build without the evidence directory says what is absent.
\newcommand{\pendingfigure}[2]{%
  \begin{center}\setlength{\fboxsep}{6pt}%
  \fbox{\parbox{.92\linewidth}{\centering\small Pending figure. #1\par\smallskip
  \footnotesize Data: \texttt{\detokenize{#2}}}}\end{center}}
\newcommand{\datafigure}[3]{\IfFileExists{#1}{#2}{\pendingfigure{#3}{#1}}}
% ---------------------------------------------------------------------------
% Terminology. paper/terminology.tex declares every term once with
% \term{<scope>}{<term>}{<definition>}, scope main (terms the main text relies on) or app
% (terms used only in the appendices); the glossary (Appendix app:glossary) prints both
% scopes. \termrows{<scope>} expands to the
% rows "term & definition \\" of that scope, for use inside a two-column tabular.
\makeatletter
\@namedef{termrows@main}{}
\@namedef{termrows@app}{}
\newcommand{\term}[3]{\expandafter\g@addto@macro\csname termrows@#1\endcsname{#2 & #3\\}}
\newcommand{\termrows}[1]{\@nameuse{termrows@#1}}
\makeatother
